\section{はじめに}
画像データセットに対する物体アノテーションは，物体が画像のどこに存在するかという物体領域の判断と，その物体が何であるかというカテゴリの判断から構成される．
YOLO\cite{ref_YOLO}やDETR\cite{ref_DETR}に代表される物体検出モデルは，両者を学習時に定めた固定のクラス集合に対して行うため，COCO\cite{ref_COCO}の80クラスで学習したモデルは，80クラスより細かい語や80クラスの外の語を出力できない．
新しいクラスを扱うには，そのクラスの注釈を集めて再学習する必要がある．

著者らの前稿\cite{ref_prev,ref_prevVG}では，物体領域の検出とカテゴリのアノテーションを分離し，検出した領域に視覚言語モデル（Vision-Language Model，VLM）のゼロショット分類を適用して，WordNet\cite{ref_WordNet}の下位語をたどることで再学習なしにカテゴリの粒度を変更する手法を検討した．
前稿の構成は，DETRのボックスの切り抜き画像にCLIP\cite{ref_CLIP}を適用するものであった．
本稿では，検出器をOWLv2\cite{ref_OWLv2}のクラス非依存の物体らしさに，分類器を注目領域をマスクで与えるAlpha-CLIP\cite{ref_AlphaCLIP}に置き換え，前稿\cite{ref_prevVG}の\rev{WordNetとVisual Genome（VG）\cite{ref_VG}を組み合わせた語彙の木}で粒度を決めるモデルを示す．

このモデルでは，検出器がアノテーションできる物体の範囲を決め，分類器と語彙の木が出力する語の正しさと細かさを決める．
COCOで学習したYOLOやDETRはCOCOの物体を少ないボックスで正確に囲むが，COCO外の物体の領域を出力するかは明らかでない．
OWLv2の物体らしさは学習時のクラスに依存しないが，ボックスの数が多くなる．
また，CLIPに切り抜き画像を入力すると周囲の文脈が失われ，Alpha-CLIPで領域を指定する場合との差は明らかでない．
そこで本稿では，モデルの各部品を置き換えた構成と比較し，部品を選ぶ理由を実験により示す．

本稿の貢献は以下の3点である．
第一に，提案手法の検出器をOWLv2，YOLO11x，DETRの3つで比較し，COCOの物体の再現率，COCO外の物体の再現率，検出から命名までのF値を示す．
これにより，COCOの物体に限ればYOLO11xとDETRのボックスが有利である一方，COCO外の物体を囲めるのはOWLv2のみであることを示し，OWLv2を選ぶ理由を明らかにする．
第二に，分類器をCLIPとAlpha-CLIPで比較し，いずれの検出器でもAlpha-CLIPのF値と，COCOのクラスより細かい語を正しく出力する割合が高いことを示す．
第三に，\rev{WordNet$ \cap $VGの木（\ref{sec_vocab}節）}の部分木ごとに語の確率を合計して出力の深さを決める命名が，WordNetの下位語を1段ずつたどる命名より細かい粒度の評価を高めることを，GTのボックスで示す．
なお，検出したボックスのうちどれを残すかを決める棄却の判定について，本稿では単純な規則を用い，その設計は今後の課題とする．
